\subsection{Comprobación de la fecha objetivo de E2}
La revisión distingue el cierre de evidencia, el pronunciamiento contractual y la ejecución del corte. Para inicio el 1 de diciembre de 2026, la evidencia de las cuatro últimas semanas de marcha blanca termina el 31 de julio de 2028 y la producción/Operación simultáneas tienen como objetivo el 1 de agosto. El artículo 18.3 reconoce hasta diez días hábiles de revisión y otros diez de subsanación; son ventanas máximas, no un mínimo obligatorio de espera. La preparación diaria puede reducir trabajo de revisión, pero no impone al CLIENTE un pronunciamiento inmediato sobre evidencia aún no producida.\parencite[arts.~17.2.4, 17.3 y 18.3]{bases-admin}

\begin{tabularx}{\linewidth}{Y L{27mm} L{27mm}}\hline
Escenario & Fin revisión & Fin subsanación\\\hline
Sin observaciones, revisión máxima & 14/08/2028 & No aplica\\\hline
Con observaciones, ambas ventanas máximas & 14/08/2028 & 29/08/2028\\\hline
\end{tabularx}\par\smallskip
El cálculo usa lunes a viernes y feriados chilenos del escenario; sus datos se conservan en \path{componentes/recepcion_e2_contraste_contractual.csv}. No incluye una nueva revisión ni la duración efectiva del corte. Incluso sin observaciones, el uso completo del plazo de revisión termina después de la fecha objetivo. Agregar recursos a Synaptix, adelantar informes parciales o reducir una reserva no elimina esa dependencia externa.

\paragraph{Resultado y responsabilidad.} El escenario conservador no satisface el objetivo del primer día de 21; el corte del 31 de agosto del contraste de revisión posterior es un escenario de atraso, no la fecha contractual ofrecida. La red seleccionada programa acta de marcha blanca y corte el 1 de agosto; H12 es la aceptación final de implementación dentro de 21. La línea base adopta el supuesto de participación de la Contraparte en la recepción continua, la sesión final del 31 de julio y la firma expresa al cierre de la evidencia, antes del corte. Así se cierra la selección del procedimiento y del calendario ofrecido; la firma y la aceptación son actuaciones futuras sujetas a los criterios de las bases. Proyecto coordinará la ejecución y PL-R04 controlará cualquier desviación. El contraste de ventanas máximas cuantifica el atraso si el supuesto no se cumple; no sustituye el escenario seleccionado ni reduce derechos de revisión. Los servicios de ambos alcances se dimensionan desde el primer día de 21; su dotación no acredita por sí sola que E2 haya sido recibido. No se sustituyen las últimas mediciones por proyecciones ni se desplazan los meses contractuales.\parencite[arts.~17--18]{bases-admin}
